% Draft body for the competition's ACM sigconf Overleaf template (overleaf.com/read/vjztpvrnbrdh).
% Paste the sections into the template; keep its fixed title/subtitle/intro/acknowledgment text.
% Limit: 4 pages excluding references. \todo{} marks numbers to update for the final submission.
% Sources: docs/EXPERIMENTS.md (OOF), runs/jobs/submissions.csv (LB), docs/COMPETITION_NOTES.md.
\newcommand{\todo}[1]{\textcolor{red}{#1}}

\begin{abstract}
We segment solar filaments in full-disk GONG H$\alpha$ images in two stages: dense U-Net probability
maps at native resolution, then learned instance selection and per-instance refinement. Stage 1
averages ConvNeXt U-Nets trained on native-resolution crops ($512^2$ and $768^2$) with full-disk
$1024^2$ models, all trained on soft annotator-consensus targets. Stage 2 builds a tree of candidate
instances from 10 threshold levels, scores each with a gradient-boosted model and a crop refiner aligned
with the filament's principal axis, and keeps the subset that maximises expected panoptic quality. With
5-fold time-grouped cross-validation the accuracy configuration (8 stage-1 members) reaches an
out-of-fold PQ of 0.467 at 18.7\,s per image and the efficient configuration (3 members) 0.462 at
3.3\,s per image on one GPU (one expert scored against the other experts: 0.336). Both score 0.40 on
the public leaderboard, using only the supplied test image.
\end{abstract}

\section{Data, metric and what limits the score}
The training set has 707 frames from six GONG sites (2011--2022) with 1,154 annotation sets by 37
annotators; 58\% of the frames have a single annotator. PQ is computed against every annotation set
separately and pooled, so a prediction is rewarded for matching each annotator's filament with
IoU $>0.5$. Two properties of the labels shape the design.
\textbf{(1) Annotators disagree.} One expert scored against the others reaches PQ 0.336 (SQ 0.634,
RQ 0.530). Annotators draw 4.3 to 10.5 filaments per frame on average, their median filament area
differs by $2\times$, and the annotator of a frame cannot be inferred from its site or year.
We therefore predict the \emph{consensus}: targets are the per-pixel mean over annotators.
\textbf{(2) The annotation protocol} \cite{magfilo} skips filaments beyond $\pm70^\circ$ from the
central meridian, small grouped filaments and too-faint ones, and asks for tight, single-piece,
hole-free outlines that include barbs. In the training labels only 0.05\% of filaments lie beyond
$70^\circ$ longitude. Whether a fragmented filament is one instance or several is annotator-dependent
and is the main source of merge/split errors.

\section{Method}
\begin{figure*}[t]
\includegraphics[width=\textwidth]{fig1_pipeline.pdf}
\caption{Top: one out-of-fold frame through the pipeline (native-resolution crop around its largest filament).
(c) stage-1 ensemble probability; (d) the candidate tree is the nested set of connected components over
10 thresholds (three shown); (e) the selected and refined filament (green) against the outlines of the three
annotators. The candidate on the right, which no annotator drew, is rejected by the scorer.
Bottom: pipeline stages. Produced by \texttt{tools/make\_figure\_pipeline.py}.}
\label{fig:pipeline}
\end{figure*}

\paragraph{Pre-processing.} Each frame is disk-fitted, limb-darkening corrected with a radial profile
and contrast-normalised. Only the supplied image is used.

\paragraph{Stage 1: probability maps.} U-Nets with three heads (filament mask, outline band, spine),
trained with BCE + soft Dice on soft consensus targets. The main members see native-resolution crops
($512^2$ or $768^2$ of the $2048^2$ frame; Hann-blended sliding-window inference); full-disk $1024^2$
members add global context and remove false positives. Encoders are ImageNet-22k ConvNeXt.
We evaluated two configurations; A is our final submission and B is its efficient variant
(also scored on the public leaderboard). \textbf{A (accuracy):} 8 members $\times$ 5 folds (three seeds of
ConvNeXt-T and one ConvNeXt-S on $512^2$ crops, ConvNeXt-T and ConvNeXt-B on $768^2$ crops, two full-disk
ConvNeXt-T seeds), with the 8 dihedral transforms as test-time augmentation.
\textbf{B (efficiency):} 3 members $\times$ 5 folds (two $512^2$-crop ConvNeXt-T seeds, one full-disk
ConvNeXt-T) with horizontal-flip augmentation; its scorer and $\tau$ are tuned on flip-only OOF maps.

\paragraph{Stage 2: candidate tree and learned selection.} Thresholding the ensemble map at 10 levels
(0.1--0.9) gives nested connected components, about 100 candidates per frame. A histogram GBM predicts
each candidate's mean over annotators of IoU$\cdot[\mathrm{IoU}>0.5]$ from 42 features (probability
statistics, shape, boundary sharpness, image contrast, disk position, tree context). Keeping a candidate
with expected contribution $q$ raises PQ iff $q>\mathrm{PQ}/2$; an exact dynamic program on the tree keeps
at most one candidate per root-to-leaf path with $q>\tau$.

\paragraph{Rotated crop refiner.} Each candidate is re-segmented by a ResNet34 U-Net on a
$384\times160$ crop aligned with its principal axis (inputs: image, stage-1 probability, candidate mask),
which predicts the consensus outline of \emph{this} filament and a score. Axis-aligned crops spend most
pixels on background for thin, tilted filaments. The final score blends the GBM and refiner scores;
refined masks are painted in score order without overlap.

\section{Validation protocol}
All stages are validated out-of-fold on 5 folds grouped by 14-day windows, so a filament never appears in
both training and validation. Downstream stages (scorer, refiner) are trained only on OOF maps. PQ is
computed with a re-implementation tested against the organisers' notebook. The OOF$\rightarrow$public offset
is a stable $-0.04$ to $-0.07$ (the public score is rounded to two decimals), and every change that improved OOF improved or tied the public score
(Table~\ref{tab:ablation}).

\section{Results}
\begin{table}[t]
\caption{Ablation. OOF PQ over 707 frames; public leaderboard on $\sim$90 test frames (2 decimals).}
\label{tab:ablation}
\small
\begin{tabular}{lcc}
\toprule
Pipeline & OOF PQ & Public \\
\midrule
ResNet34, full disk $1024^2$, threshold & 0.363 & 0.31 \\
+ learned candidate selection (GBM) & 0.378 & 0.32 \\
ConvNeXt-T, $512^2$ native crops, threshold & 0.425 & 0.37 \\
Ensemble (crops + full disk), threshold & 0.445 & 0.38 \\
+ second crop-model seed, threshold & 0.449 & 0.39 \\
+ GBM selection + rotated refiner & 0.465 & 0.40 \\
+ $768^2$ crops, ConvNeXt-B, more seeds (8 members) & 0.467 & 0.40 \\
\bottomrule
\end{tabular}
\end{table}

\paragraph{What helped.} Native resolution (+0.03 over $1024^2$ full disk), ConvNeXt-T over ResNet34
(+0.017 on fold 0), a full-disk member against false positives, a second crop-model seed (+0.004),
learned selection instead of one threshold, and refining each instance in a rotated frame (+0.005).
$768^2$ crops gain +0.007 as a single model but nothing once blended with a full-disk member, which
already supplies the context.
\paragraph{What did not help.} An object-level BiCC loss (only shifted calibration), larger encoders
(ConvNeXt-S/B: within seed noise once blended), a ResNet50 refiner, grouping fragments before selection,
disk-longitude features (the GBM already used disk position), and sub-pixel target offsets.
\paragraph{Test-time choices checked on the public set.} Moving $\tau$ from the OOF optimum to 0.22 or
0.38, or dropping the refiner, each cost 0.01 on the public set, so the OOF optimum transfers.
Horizontal-flip TTA instead of the 8 dihedral views tied (0.40) at half the cost.

\section{Morphology and errors}
\begin{figure}[t]
% files: runs/figures/v11_refined/oof_20110429135554Bh_zoom1.jpg, oof_20110128174814Mh_zoom0.jpg
\includegraphics[width=0.49\columnwidth]{oof_20110429135554Bh_zoom1.jpg}\hfill
\includegraphics[width=0.49\columnwidth]{oof_20110128174814Mh_zoom0.jpg}
\caption{Native-resolution out-of-fold crops (configuration A). The three annotators' outlines (yellow, cyan,
magenta) disagree on barbs and ends; the refined prediction (green = matched) follows the consensus body.}
\label{fig:morphology}
\end{figure}
On the final OOF predictions of configuration A (organisers' metric, 1,154 annotation sets) the system reaches a mean
foreground Dice of 0.669 per annotation set (one expert against the others: 0.612), a median matched IoU
of 0.69 (expert: 0.63), and 64.7\% of GT filaments matched with IoU $>0.5$ (21.6\% untouched, 3.8\%
with IoU $<0.3$, 10.0\% near misses at 0.3--0.5). Fragmentation is rare: 2.5\% of GT filaments touch
more than one prediction and 1.9\% of predictions touch more than one GT filament (expert: 3.7\% / 2.6\%).
Missed filaments are mostly small ($<700$\,px) or drawn by only one of several annotators: in an earlier
error analysis of the ensemble, filaments drawn by all annotators were matched 81\% of the time, those
drawn by one of several 18\%.

\section{Efficiency}
Test inference on one RTX 4090, uncached, per $2048^2$ image:
\begin{tabular}{lccc}
\toprule
Configuration & s/image & OOF PQ & Public \\
\midrule
8 stage-1 models $\times$ 5 folds, D4 TTA & 18.7 & 0.467 & 0.40 \\
8 models $\times$ 5 folds, hflip TTA & 6.7 & -- & 0.40 \\
3 models $\times$ 5 folds, hflip TTA (tuned on hflip OOF maps) & 3.3 & 0.462 & 0.40 \\
\bottomrule
\end{tabular}
Stage 1 dominates the cost: per image, the fast configuration spends 0.2\,s on pre-processing, 1.4\,s on
stage 1 (two 512-crop ConvNeXt-T members 0.65\,s each, the full-disk member 0.1\,s), 0.1\,s on candidates and
GBM, and 0.8\,s on the five refiners; the rest is image I/O and RLE encoding. D4 TTA costs 3.6$\times$ hflip
for the crop members. A paired bootstrap over the 14-day groups of the OOF set puts A $-$ B at
$+0.0045$ PQ (95\% interval $+0.0016$ to $+0.0073$; A is ahead on all five folds), so the 5.7$\times$
cost of A buys a small but real gain that the two-decimal public score cannot resolve. Refiner batches are padded to a fixed size, so cuDNN autotuning runs once (it had cost
25\,s per image).

\section{Reproducibility and integrity}
Code (MIT): \todo{repo URL}; a notebook reproduces the submission from the released checkpoints.
The public MAGFiLO release contains labels for 139 of the 180 test images. On the training set, copying
one released annotation set per frame scores PQ 0.757, which explains public scores near 0.70.
We used neither these labels nor any ground truth beyond the training set.
